\documentclass[11pt]{article}

\usepackage{amsmath, amssymb, bm}
\usepackage{geometry}
\geometry{margin=2.5cm}

\title{Mathematical Formulation of the Restricted Hartree--Fock Algorithm Implemented in \texttt{QuetzalQC}}
\author{}
\date{}

\begin{document}

\maketitle

\section{Electronic Structure Problem}

Within the Born--Oppenheimer approximation, the electronic Hamiltonian for a fixed nuclear configuration is

\begin{equation}
\hat{H}_{elec} =
\sum_{i}^{N} \left(
-\frac{1}{2} \nabla_i^2
- \sum_{A} \frac{Z_A}{|\mathbf{r}_i - \mathbf{R}_A|}
\right)
+
\sum_{i<j}^{N} \frac{1}{|\mathbf{r}_i - \mathbf{r}_j|}.
\end{equation}

The electronic Schrödinger equation is

\begin{equation}
\hat{H}_{elec} \Psi = E \Psi.
\end{equation}

An exact solution is intractable for multi-electron systems, therefore approximations are required.

\section{Hartree--Fock Approximation}

In the Restricted Hartree--Fock (RHF) approach, the many-electron wavefunction is approximated by a single Slater determinant constructed from spatial orbitals:

\begin{equation}
\Psi = \frac{1}{\sqrt{N!}} \det \left[ \chi_i(\mathbf{x}_j) \right],
\end{equation}

where spin-orbitals are products of spatial orbitals and spin functions.

For closed-shell systems, each spatial orbital is doubly occupied.

\section{Basis Set Expansion}

Each molecular orbital is expanded in a finite atomic orbital (AO) basis:

\begin{equation}
\phi_i(\mathbf{r}) = \sum_{\mu}^{K} C_{\mu i} \, \chi_\mu(\mathbf{r}),
\end{equation}

where $\{ \chi_\mu \}$ are basis functions (e.g., STO-3G).

\section{Roothaan--Hall Equations}

The Hartree--Fock equations in a non-orthogonal AO basis reduce to the generalized eigenvalue problem:

\begin{equation}
\mathbf{F} \mathbf{C} = \mathbf{S} \mathbf{C} \bm{\varepsilon},
\end{equation}

where

\begin{itemize}
\item $\mathbf{S}_{\mu\nu} = \langle \chi_\mu | \chi_\nu \rangle$ is the overlap matrix,
\item $\mathbf{F}$ is the Fock matrix,
\item $\bm{\varepsilon}$ is the diagonal matrix of orbital energies.
\end{itemize}

\section{Fock Matrix Construction}

The Fock matrix is defined as

\begin{equation}
F_{\mu\nu} = H_{\mu\nu} + J_{\mu\nu} - \frac{1}{2} K_{\mu\nu},
\end{equation}

where

\begin{equation}
H_{\mu\nu} = T_{\mu\nu} + V_{\mu\nu}
\end{equation}

is the one-electron Hamiltonian matrix.

The Coulomb and exchange matrices are given by

\begin{equation}
J_{\mu\nu} =
\sum_{\lambda\sigma}
P_{\lambda\sigma}
(\mu\nu|\lambda\sigma),
\end{equation}

\begin{equation}
K_{\mu\nu} =
\sum_{\lambda\sigma}
P_{\lambda\sigma}
(\mu\lambda|\nu\sigma),
\end{equation}

where $(\mu\nu|\lambda\sigma)$ are the two-electron repulsion integrals.

\section{Density Matrix}

For a closed-shell system, the density matrix is

\begin{equation}
P_{\mu\nu} =
2 \sum_{i}^{occ}
C_{\mu i} C_{\nu i}.
\end{equation}

\section{Self-Consistent Field Procedure}

The algorithm iteratively solves:

\begin{enumerate}
\item Construct $\mathbf{F}$ from current $\mathbf{P}$.
\item Solve $\mathbf{F}\mathbf{C} = \mathbf{S}\mathbf{C}\bm{\varepsilon}$.
\item Build new density matrix $\mathbf{P}$.
\item Repeat until convergence:
\[
\|\mathbf{P}^{(n+1)} - \mathbf{P}^{(n)}\| < \delta.
\]
\end{enumerate}

This defines a fixed-point problem:

\begin{equation}
\mathbf{P} = \mathcal{F}(\mathbf{P}).
\end{equation}

\section{Total Energy Expression}

The RHF total energy is

\begin{equation}
E =
\frac{1}{2}
\mathrm{Tr}
\left[
\mathbf{P}
(\mathbf{H} + \mathbf{F})
\right]
+ E_{nuc},
\end{equation}

where $E_{nuc}$ is the nuclear repulsion energy.

\section{Computational Scaling}

The dominant cost arises from evaluation of the Coulomb and exchange contractions:

\[
\mathcal{O}(K^4),
\]

where $K$ is the number of basis functions.

\section{Current Scope}

The present implementation:

\begin{itemize}
\item Uses STO-3G basis.
\item Solves the RHF equations for small molecules (e.g., H$_2$).
\item Employs a modular SCF architecture.
\item Validates results against reference implementations.
\end{itemize}

\end{document}

